% "{'chapitre':'slci_multiphy','classe':('PSI'),'type':('application'),'titre':'Mise à l\\'eau d\\'un robot sous-marin','source':'Concours Centrale - MP 2019','comp':['B2-02',],'corrige':True}"

%\setchapterimage{fig_00.jpg}
\chapter*{Application \arabic{cptApplication} \\ 
Prothèse active transtibiale -- \ifprof Corrigé \else Sujet \fi}
\addcontentsline{toc}{section}{Application \arabic{cptApplication} : Prothèse active transtibiale-- \ifprof Corrigé \else Sujet \fi}

\iflivret \stepcounter{cptApplication} \else
\ifprof  \stepcounter{cptApplication} \else \fi
\fi

\setcounter{question}{0}
\marginnote{Concours CCMP -- PSI}





\subsection*{Question de cours}
\question{Donner la strucutre de la chaîne de puissance. On précisera en particulier le nom de chacune des fonctions.}

\question{Donner 3 exemples de composants permettant de réaliser chacune des fonctions.}

\subsection*{Prothèse active transtibiale}
\marginnote{D'après ressources de Florestan Mathurin.}

\begin{marginfigure}
\includegraphics[width=\linewidth]{fig_01}
\end{marginfigure}

La majorité des prothèses transtibiales (pour une amputation en
dessous du genou) utilisées aujourd'hui sont purement passives, c'està-
dire que leurs propriétés mécaniques restent fixes pendant la
marche. Ces prothèses sont constituées en général de semelles ressorts
en carbone profilées qui emmagasinent et restituent l'énergie
mécanique pendant la marche par déformation.
On s’intéresse ici à un prototype mis au point par des ingénieurs du
MIT qui a permis la mise au point d’une nouvelle génération de
prothèse, dite active. Cette prothèse active transtibiale est capable de
proposer un comportement similaire à celui des membres non
amputés.

L'actionneur de la prothèse est un moteur à courant
continu alimenté par une batterie rechargeable de 16
Volts. L'énergie mécanique est transmise par un
réducteur de type poulies-courroie suivi d'un système
vis-écrou qui adapte cette énergie mécanique pour la
prothèse (ensemble de liaisons entre le pied artificiel
constitué d'une semelle en fibres de carbone et le
manchon ou tibia artificiel). Des ressorts permettent
d'ajuster également l'énergie mécanique fournie au
pied artificiel. L'effort exercé par les ressorts est
directement relié au couple exercé par l'actionneur.



\begin{center}
\includegraphics[width=.45\linewidth]{fig_02}
\includegraphics[width=.45\linewidth]{fig_03}
\end{center}

Les informations délivrées par les capteurs sont
traitées par un calculateur qui élabore la commande
du moteur.


La chaîne d'informations est constituée d'un ensemble de capteurs permettant d'acquérir différentes informations :
\begin{itemize}
\item un potentiomètre linéaire qui mesure l'allongement/écrasement des ressorts;
\item un codeur incrémental placé au niveau de l'articulation pied/tibia;
\item plusieurs capteurs capacitifs disposés sous la semelle du pied au niveau du talon (2 capteurs) et à l'avant du pied (4 capteurs).
\end{itemize}
\question{Compléter le diagramme chaîne d’information / chaîne de puissance.}

\begin{center}
\includegraphics[width=\linewidth]{fig_04}
\end{center}


